%
% einleitung.tex -- Beispiel-File für die Einleitung
%
% (c) 2020 Prof Dr Andreas Müller, Hochschule Rapperswil
%
% !TEX root = ../../buch.tex
% !TEX encoding = UTF-8
%
\section{Das dritte keplersche Gesetz
  \label{kepler:section:einleitung}}

Das dritte keplersche Gesetz beschreibt den Zusammenhang zwischen der Umlaufzeit eines Planeten und der Grösse seiner Umlaufbahn \cite{vogt1995}.
Die Regel besagt, dass sich die Quadrate der Umlaufzeiten jeglicher Planeten zueinander verhalten wie die Kuben ihrer mittleren Abstände zur Sonne \cite{vogt1995}.
Mathematisch lässt sich dies für einen Planeten mit der Umlaufzeit \(T\) und der grossen Halbachse \(a\) als
\begin{equation}
    T^2 = \frac{4\pi^2}{GM} a^3
    \label{kepler:equation:drittes_gesetz}
\end{equation}
ausdrücken \cite{vogt1995}.
Die Konstante auf der rechten Seite der Gleichung \eqref{kepler:equation:drittes_gesetz} hängt massgeblich von der Gravitationskonstante \(G\) und der Masse des Zentralgestirns \(M\) ab \cite{vogt1995}.
Im Folgenden werden verschiedene mathematische Wege aufgezeigt, um diese fundamentale physikalische Beziehung herzuleiten.